% GENERATED FILE. Edit canonical lessons, then run npm run generate:books.
% canonical-source-hash: fnv1a64:1f77fc6ceb91b04f
% canonical-lessons: HI-C408-mil, HI-C408-kos, HI-C408-fasla, HI-C408-rah, HI-C408-bhagana

\chapter{Miles and Distances}
\label{ch:hi-miles-and-distances}
\hlchaptermodality{408}

\begin{chapteropening}
\textbf{By the end of this chapter:} \emph{I can say how far apart two places are, in miles or in kos, and say a driver went too fast.}

Complete the last lesson of chapter 408: I can say how far apart two places are, in miles or in kos, and say a driver went too fast.
\end{chapteropening}

\section[\texorpdfstring{\hi{मील} (mīl)}{मील (mīl)}]{\hi{मील} (mīl) — a mile}

\label{lesson:HI-C408-mil}



\begin{quote}
Before the new one: say the Hindi for a convoy, then the Hindi for to wash away.
\end{quote}

\subsection*{You'll want to know: \hi{मील}}
\textbf{\hi{मील}} — \emph{mīl} — \textquotedblleft{}a mile\textquotedblright{}.

\begin{quote}
\hi{स्टेशन} \hi{यहाँ} \hi{से} \hi{दो} \hi{मील} \hi{है।}
\end{quote}

\begin{quote}
\emph{sṭeśan yahā̃ se do mīl hai.}
\end{quote}

\begin{quote}
The station is two miles from here.
\end{quote}

Masculine.

\begin{cousinweb}
\textbf{\hi{मील}} (\emph{mīl}) is borrowed from English \emph{mile}; road signs now give kilometres, but people still count in miles.
\end{cousinweb}

\begin{grammarlens}[title={What you've built}]
One more word for getting around.
\end{grammarlens}

\subsection*{Your turn}
\begin{itemize}
\raggedright
  \item \emph{Say it:} \emph{mīl}
  \item \emph{Say it:} \emph{mīl}, once more
  \item \emph{Say it:} \emph{sṭeśan yahā̃ se do mīl hai.} — The station is two miles from here.
  \item \emph{Recall:} say the Hindi for a convoy, then the Hindi for to wash away, then say \emph{mīl} again
\end{itemize}

\begin{tcolorbox}[breakable,colback=teal!4,colframe=teal!35!black,title={Before you move on}]
What does \hi{मील} mean? (\textquotedblleft{}A mile\textquotedblright{}.) Say it once more.
\end{tcolorbox}
\section[\texorpdfstring{\hi{कोस} (kos)}{कोस (kos)}]{\hi{कोस} (kos) — a kos (about two miles)}

\label{lesson:HI-C408-kos}



\begin{quote}
Before the new one: say the Hindi for to wash away, then the Hindi for a mile.
\end{quote}

\subsection*{You'll want to know: \hi{कोस}}
\textbf{\hi{कोस}} — \emph{kos} — \textquotedblleft{}a kos (about two miles)\textquotedblright{}.

\begin{quote}
\hi{दादा} \hi{कहते} \hi{थे} \hi{कि} \hi{नदी} \hi{तीन} \hi{कोस} \hi{दूर} \hi{है।}
\end{quote}

\begin{quote}
\emph{dādā kahte the ki nadī tīn kos dūr hai.}
\end{quote}

\begin{quote}
Grandfather used to say that the river is three kos away.
\end{quote}

Masculine.

\begin{cousinweb}
\textbf{\hi{कोस}} (\emph{kos}) is a tadbhava of Sanskrit \emph{krośa}, \textquotedblleft{}a shout\textquotedblright{}, and then \textquotedblleft{}the distance a shout carries\textquotedblright{} (CDIAL 3611), which lists Hindi \emph{kos}, about two miles.
\end{cousinweb}

\begin{grammarlens}[title={What you've built}]
One more word for getting around.
\end{grammarlens}

\subsection*{Your turn}
\begin{itemize}
\raggedright
  \item \emph{Say it:} \emph{kos}
  \item \emph{Say it:} \emph{kos}, once more
  \item \emph{Say it:} \emph{dādā kahte the ki nadī tīn kos dūr hai.} — Grandfather used to say that the river is three kos away.
  \item \emph{Recall:} say the Hindi for to wash away, then the Hindi for a mile, then say \emph{kos} again
\end{itemize}

\begin{tcolorbox}[breakable,colback=teal!4,colframe=teal!35!black,title={Before you move on}]
What does \hi{कोस} mean? (\textquotedblleft{}A kos (about two miles)\textquotedblright{}.) Say it once more.
\end{tcolorbox}
\section[\texorpdfstring{\hi{फ़ासला} (fāslā)}{फ़ासला (fāslā)}]{\hi{फ़ासला} (fāslā) — a distance; a gap}

\label{lesson:HI-C408-fasla}



\begin{quote}
Before the new one: say the Hindi for a mile, then the Hindi for a kos.
\end{quote}

\subsection*{You'll want to know: \hi{फ़ासला}}
\textbf{\hi{फ़ासला}} — \emph{fāslā} — \textquotedblleft{}a distance; a gap\textquotedblright{}.

\begin{quote}
\hi{दोनों} \hi{गाँवों} \hi{के} \hi{बीच} \hi{दस} \hi{मील} \hi{का} \hi{फ़ासला} \hi{है।}
\end{quote}

\begin{quote}
\emph{donõ gā̃voṁ ke bīc das mīl kā fāslā hai.}
\end{quote}

\begin{quote}
The two villages are ten miles apart.
\end{quote}

Masculine.

\begin{cousinweb}
\textbf{\hi{फ़ासला}} (\emph{fāslā}) is usually traced to Arabic \emph{fāṣila}, \textquotedblleft{}a separating space, an interval\textquotedblright{}, which came to Hindi through Persian.
\end{cousinweb}

\begin{grammarlens}[title={What you've built}]
One more word for getting around.
\end{grammarlens}

\subsection*{Your turn}
\begin{itemize}
\raggedright
  \item \emph{Say it:} \emph{fāslā}
  \item \emph{Say it:} \emph{fāslā}, once more
  \item \emph{Say it:} \emph{donõ gā̃voṁ ke bīc das mīl kā fāslā hai.} — The two villages are ten miles apart.
  \item \emph{Recall:} say the Hindi for a mile, then the Hindi for a kos, then say \emph{fāslā} again
\end{itemize}

\begin{tcolorbox}[breakable,colback=teal!4,colframe=teal!35!black,title={Before you move on}]
What does \hi{फ़ासला} mean? (\textquotedblleft{}A distance; a gap\textquotedblright{}.) Say it once more.
\end{tcolorbox}
\section[\texorpdfstring{\hi{राह} (rāh)}{राह (rāh)}]{\hi{राह} (rāh) — a path; a way (literary)}

\label{lesson:HI-C408-rah}



\begin{quote}
Before the new one: say the Hindi for a kos, then the Hindi for a distance.
\end{quote}

\subsection*{You'll want to know: \hi{राह}}
\textbf{\hi{राह}} — \emph{rāh} — \textquotedblleft{}a path; a way (literary)\textquotedblright{}.

\begin{quote}
\hi{यह} \hi{राह} \hi{जंगल} \hi{से} \hi{होकर} \hi{जाती} \hi{है।}
\end{quote}

\begin{quote}
\emph{yah rāh jaṅgal se hokar jātī hai.}
\end{quote}

\begin{quote}
This path goes through the forest.
\end{quote}

Feminine.

Songs and sayings prefer \hi{राह}; in plain talk people say \hi{रास्ता}.

\begin{cousinweb}
\textbf{\hi{राह}} (\emph{rāh}) is usually traced to Persian \emph{rāh}, \textquotedblleft{}a road, a way\textquotedblright{}, the word inside \hi{चौराहा}, where four ways meet.
\end{cousinweb}

\begin{grammarlens}[title={What you've built}]
One more word for getting around.
\end{grammarlens}

\subsection*{Your turn}
\begin{itemize}
\raggedright
  \item \emph{Say it:} \emph{rāh}
  \item \emph{Say it:} \emph{rāh}, once more
  \item \emph{Say it:} \emph{yah rāh jaṅgal se hokar jātī hai.} — This path goes through the forest.
  \item \emph{Recall:} say the Hindi for a kos, then the Hindi for a distance, then say \emph{rāh} again
\end{itemize}

\begin{tcolorbox}[breakable,colback=teal!4,colframe=teal!35!black,title={Before you move on}]
What does \hi{राह} mean? (\textquotedblleft{}A path; a way (literary)\textquotedblright{}.) Say it once more.
\end{tcolorbox}
\section[\texorpdfstring{\hi{भगाना} (bhagānā)}{भगाना (bhagānā)}]{\hi{भगाना} (bhagānā) — to chase off; to drive (a vehicle) fast}

\label{lesson:HI-C408-bhagana}



\begin{quote}
Before the new one: say the Hindi for a distance, then the Hindi for a path.
\end{quote}

\subsection*{You'll want to know: \hi{भगाना}}
\textbf{\hi{भगाना}} — \emph{bhagānā} — \textquotedblleft{}to chase off; to drive (a vehicle) fast\textquotedblright{}.

\begin{quote}
\hi{ड्राइवर} \hi{ने} \hi{गाड़ी} \hi{बहुत} \hi{तेज़} \hi{भगाई।}
\end{quote}

\begin{quote}
\emph{ḍrāivar ne gāṛī bahut tez bhagāī.}
\end{quote}

\begin{quote}
The driver drove the car very fast.
\end{quote}

Past, with \hi{ने} (he or she …): \hi{उसने} \hi{कुत्ते} \hi{को} \hi{भगाया}, \emph{usne kutte ko bhagāyā}.

\begin{cousinweb}
\textbf{\hi{भगाना}} (\emph{bhagānā}) is the causative partner of \hi{भागना}, \textquotedblleft{}to run away\textquotedblright{} (CDIAL 9361): to make a dog run off, or to make a car run fast.
\end{cousinweb}

\begin{grammarlens}[title={What you've built}]
One more word for getting around.
\end{grammarlens}

\subsection*{Your turn}
\begin{itemize}
\raggedright
  \item \emph{Say it:} \emph{bhagānā}
  \item \emph{Say it:} \emph{bhagānā}, once more
  \item \emph{Say it:} \emph{ḍrāivar ne gāṛī bahut tez bhagāī.} — The driver drove the car very fast.
  \item \emph{Recall:} say the Hindi for a distance, then the Hindi for a path, then say \emph{bhagānā} again
\end{itemize}

\begin{tcolorbox}[breakable,colback=teal!4,colframe=teal!35!black,title={Before you move on}]
What does \hi{भगाना} mean? (\textquotedblleft{}To chase off; to drive (a vehicle) fast\textquotedblright{}.) Say it once more.
\end{tcolorbox}
